% ============================================================
%  Generalized DM-CM Decomposition — Student Tutorial
%  Source: Brovont, IEEE Trans. Power Electron., Vol. 33, No. 8, Aug. 2018
%  Figures stored in: figures/
%  Naming: GDCMDMCEAP_Fig[N]_page[XX].png
% ============================================================
\documentclass[12pt, a4paper]{article}

% ---- core packages ----
\usepackage[utf8]{inputenc}
\usepackage[T1]{fontenc}
\usepackage{amsmath, amssymb, amsthm}
\usepackage{bm}
\usepackage{siunitx}

% ---- figures ----
\usepackage{graphicx}
\graphicspath{{figures/}}
\usepackage{float}
\usepackage{caption}
\usepackage{subcaption}

% ---- layout ----
\usepackage[a4paper, top=2.5cm, bottom=2.5cm, left=3cm, right=2.8cm]{geometry}
\usepackage{parskip}
\setlength{\parskip}{6pt}
\setlength{\parindent}{0pt}

% ---- hyperlinks ----
\usepackage[colorlinks=true, linkcolor=blue!60!black, citecolor=blue!50!black]{hyperref}
\usepackage{bookmark}
\makeatletter
\pdfstringdefDisableCommands{%
    \def\vspace#1{}%
    \def\rule#1#2{}%
    \def\color#1{}%
    \let\leavevmode@ifvmode\relax%
    \def\,{ }%
}
\makeatother

% ---- styled boxes ----
\usepackage{tcolorbox}
\tcbuselibrary{skins, breakable}

% ---- tables ----
\usepackage{booktabs}
\usepackage{array}
\usepackage{longtable}

% ---- section formatting ----
\usepackage{titlesec}
\usepackage{xcolor}
\definecolor{emidark}{RGB}{10, 55, 120}
\definecolor{emiorange}{RGB}{190, 70, 5}
\definecolor{emigreen}{RGB}{8, 95, 38}
\definecolor{emired}{RGB}{170, 25, 25}
\definecolor{lightblue}{RGB}{216, 232, 252}
\definecolor{lightorange}{RGB}{255, 244, 220}
\definecolor{lightgreen}{RGB}{216, 247, 228}
\definecolor{lightred}{RGB}{255, 228, 228}
\definecolor{lightgray}{RGB}{245, 245, 245}

\titleformat{\section}{\large\bfseries\color{emidark}}{{\thesection.}}{0.5em}{}
            [\vspace{-4pt}\rule{\textwidth}{0.5pt}]
\titleformat{\subsection}{\normalsize\bfseries\color{emidark!80!black}}{{\thesubsection}}{0.5em}{}

% ---- tcolorbox styles ----
\newtcolorbox{circuitbox}[1][Circuit Description]{
    colback=lightblue, colframe=emidark, coltitle=white,
    fonttitle=\bfseries\small, title={#1},
    breakable, enhanced, left=4pt, right=4pt, top=4pt, bottom=4pt}

\newtcolorbox{statebox}[1][Key Variables]{
    colback=lightorange, colframe=emiorange, coltitle=white,
    fonttitle=\bfseries\small, title={#1},
    breakable, enhanced, left=4pt, right=4pt, top=4pt, bottom=4pt}

\newtcolorbox{eqnbox}[1][Governing Equations]{
    colback=lightgreen, colframe=emigreen, coltitle=white,
    fonttitle=\bfseries\small, title={#1},
    breakable, enhanced, left=4pt, right=4pt, top=4pt, bottom=4pt}

\newtcolorbox{physbox}[1][Physical Interpretation]{
    colback=lightgray, colframe=gray!60, coltitle=white,
    fonttitle=\bfseries\small, title={#1},
    breakable, enhanced, left=4pt, right=4pt, top=4pt, bottom=4pt}

\newtcolorbox{warningbox}[1][Watch Out]{
    colback=lightred, colframe=emired, coltitle=white,
    fonttitle=\bfseries\small, title={#1},
    breakable, enhanced, left=4pt, right=4pt, top=4pt, bottom=4pt}

% ---- custom commands ----
% CM / DM quantities
\newcommand{\vcm}{v_{\mathrm{CM}}}
\newcommand{\icm}{i_{\mathrm{CM}}}
\newcommand{\vdm}[2]{v_{#1#2}}
\newcommand{\idm}[2]{i_{#1#2}}
% Transformation matrices
\newcommand{\TiN}{\mathbf{T}^{i}_{N}}
\newcommand{\TvN}{\mathbf{T}^{v}_{N}}
% Vector notation
\newcommand{\iDCM}{\mathbf{i}_{\mathrm{DCM}}}
\newcommand{\vDCM}{\mathbf{v}_{\mathrm{DCM}}}
\newcommand{\ivec}{\mathbf{i}_{n}}
\newcommand{\vvec}{\mathbf{v}_{nP}}
% Circuit parameters
\newcommand{\Vdc}{V_{\mathrm{dc}}}
\newcommand{\Cp}{C_{p}}
\newcommand{\Rh}{R_{h}}
\newcommand{\Ni}{N_{i}}
% Figure reference
\newcommand{\figref}[1]{Figure~\ref{#1}}

% ---- figure entry macro ----
\newcommand{\figentry}[4]{%
    \begin{figure}[H]
        \centering
        \includegraphics[width=#4\textwidth]{#2}
        \caption{#3}
        \label{#1}
    \end{figure}}

% ---- document ----
\begin{document}

\begin{center}
    {\LARGE\bfseries\color{emidark}
    Differential and Common-Mode Decomposition\\[4pt]
    for Modeling Conducted EMI}\\[8pt]
    {\large A Step-by-Step Tutorial}\\[4pt]
    {\normalsize Based on: Brovont, \textit{IEEE Trans.\ Power Electron.},
    Vol.~33, No.~8, Aug.~2018}
\end{center}

\tableofcontents
\newpage

% ============================================================
\setcounter{section}{-1}
\section{Before You Begin}
% ============================================================

\subsection{What This Document Is}

Modern power electronics cannot be designed without thinking about
electromagnetic interference (EMI).  Every switch transition excites
unintended current paths through the parasitic capacitances of the
hardware, producing noise that must be suppressed to meet regulatory
limits and to protect sensitive equipment.

This tutorial takes the paper by Aaron D.\ Brovont (IEEE Transactions
on Power Electronics, August 2018) and rebuilds its logic step by
step — starting from things you already know (KVL, capacitor current,
matrix multiplication) and arriving at a systematic mathematical tool
that can predict exactly how much common-mode (CM) noise a switching
converter will produce through a parasitic baseplate capacitance.

The paper's main deliverable is a pair of \emph{transformation
matrices} that decompose any mixed-mode (MM) circuit — one where
differential-mode (DM) and CM signals are entangled — into cleanly
separated DM and CM equations.  When the circuit is asymmetric, these
equations also expose the exact \emph{coupling} between the two modes.
As a worked example, the transformation is used to build two different
(but equivalent) common-mode equivalent models for a SiC half-bridge
converter whose parasitic module capacitances are unbalanced.

By the end of this document you should be able to:
\begin{itemize}
    \item State precisely what CM and DM current and voltage mean
          mathematically.
    \item Apply the transformation matrices to a set of MM equations
          and read off the DM-to-CM coupling terms introduced by
          circuit asymmetry.
    \item Identify every element of the two CEMs in the paper and
          explain what physical mechanism each represents.
    \item Explain why the two models give identical CM currents but
          different physical pictures of the same testbed.
\end{itemize}

% ------------------------------------------------------------
\subsection{Why Switching Converters Produce EMI}
% ------------------------------------------------------------

A SiC MOSFET in a 600\,V bus can switch its drain voltage from
600\,V to 0\,V in under 100\,ns.  The rate of change is
\[
    \frac{dv}{dt} = \frac{600\,\text{V}}{100\,\text{ns}}
                  = 6 \times 10^{9}\,\text{V/s}
                  = 6\,\text{kV/}\mu\text{s}.
\]
Any capacitance $C$ connected between that switching node and the
ground reference must carry a current $i = C\,dv/dt$.  Even a
capacitance of 100\,pF — the approximate order of magnitude of a
parasitic transistor-to-baseplate capacitance in a power module —
will carry
\[
    i = 100 \times 10^{-12}\,\text{F}
        \times 6 \times 10^{9}\,\text{V/s}
      = 0.6\,\text{A}.
\]

\begin{physbox}[Why 0.6\,A Is a Problem]
This 0.6\,A is not clean sinusoidal current at the fundamental
frequency.  It is a pulse that repeats every switching period and
therefore contains harmonics at 100\,kHz, 200\,kHz, 300\,kHz, and
so on.  Regulatory conducted-emission standards (CISPR~11 for
industrial equipment) set limits as low as $-50$\,dB$\mu$A
($\approx 3\,\mu$A) at frequencies above 150\,kHz.  Three orders
of magnitude of suppression are required.  Knowing \emph{where} the
current flows is the first step to suppressing it.
\end{physbox}

This current does not flow through the intentional power path.
Instead it escapes through parasitic capacitances to the chassis, the
protective-earth wire, and the cable screen, returning through
whatever ground path is available.  This escaping current is the
\emph{common-mode conducted emission}.

% ------------------------------------------------------------
\subsection{Two Flavours of Noise: Differential Mode and Common Mode}
% ------------------------------------------------------------

Consider the simplest case: a two-wire power connection
(positive conductor and negative/return conductor) carrying current
to a load.

\begin{circuitbox}[DM and CM in a Two-Wire System]
\begin{itemize}
    \item \textbf{Differential-mode (DM) current} flows in
          \emph{opposite} directions on the two wires:
          $i_1 = -i_2$.  This is the intentional power-carrying
          current.  If you clip a current clamp around both wires
          together, DM currents cancel and the reading is zero.
          DM current is confined to the two-wire loop.

    \item \textbf{Common-mode (CM) current} flows in the \emph{same}
          direction on both wires simultaneously.  It is driven by
          the stray capacitances from both conductors to the ground
          reference.  A current clamp around both wires together
          reads \emph{twice} the CM current — the two contributions
          add instead of cancelling.  CM current escapes to ground
          and is the root cause of conducted EMI.
\end{itemize}
\end{circuitbox}

Mathematically, for a two-wire system:
\begin{align}
    i_{\mathrm{CM}} &= i_1 + i_2
        \qquad \text{(both wires, same direction)} \\
    i_{12}          &= \tfrac{1}{2}(i_1 - i_2)
        \qquad \text{(DM, push-pull component)}
\end{align}

This extends naturally to $N$ wires: the CM current is the \emph{sum}
of all line currents and the CM voltage is the \emph{average}
line-to-reference voltage.  There are $N-1$ independent DM
combinations and one CM quantity, giving $N$ quantities total — the
same count as the original $N$ line currents.  This is exactly the
transformation the paper formalises.

% ------------------------------------------------------------
\subsection{What Is a CM Equivalent Circuit and Why Do We Need One?}
% ------------------------------------------------------------

A \emph{CM equivalent circuit} (CEM) is a reduced-order model that
retains only the CM current paths of a power converter.  It replaces
the full switching-circuit detail with a handful of impedances
(capacitors, inductors, resistors) driven by voltage sources that
are functions of the switch voltages.

\begin{physbox}[The Payoff]
Once you have a CEM you can:
(1) predict CM current without running a full detailed simulation;
(2) design a CM filter to suppress it; and
(3) understand which component is the dominant noise source so you
can target design changes efficiently.
\end{physbox}

For a \emph{symmetric} system — where every component on line 1 is
identical to the corresponding component on line 2, etc.\ — deriving
a CEM is a textbook operation: parallel combinations of symmetric
impedances give the CM equivalent by inspection.  The DM and CM
equations are completely independent; you can analyse each in
isolation.

The difficulty arises when the circuit is \emph{asymmetric}.

% ------------------------------------------------------------
\subsection{The Asymmetry Problem}
% ------------------------------------------------------------

Inside a SiC multichip power module (MCPM), multiple transistor chips
share a common copper baseplate used for thermal management.  Each
chip terminal (drain, source, and phase node) has a parasitic
capacitance to that baseplate — formed by the dielectric layers in
the module package.

Due to the internal structure of the MCPM, the
\textbf{source-to-baseplate capacitance is smaller} than the
drain-to-baseplate and phase-to-baseplate capacitances.  The paper
captures this imbalance with a single asymmetry parameter $k < 1$:
the source-to-baseplate capacitance is $k\Cp$ while the other two are
each $\Cp$.

\begin{warningbox}[Why Asymmetry Is Dangerous]
In a symmetric system, DM switching waveforms drive only DM currents
and CM paths carry only CM contributions.  \textbf{Asymmetry breaks
this independence.}  A DM voltage oscillation — caused by the normal
switching action — can now ``pump'' current into the CM path.  DM
noise converts to CM noise directly through the imbalanced
capacitances.  Without a systematic model, it is impossible to predict
how much conversion occurs.
\end{warningbox}

The paper answers exactly this question: given a specific asymmetry
parameter $k$, how much DM-to-CM coupling is induced, and what does
the resulting CM equivalent circuit look like?  The answer comes in
the form of a \emph{DM-dependent voltage source} that appears inside
the CEM whenever $k \neq 1$.  When $k = 1$ (perfect symmetry) that
source equals zero and the CEM reduces to the standard symmetric
result.

% ------------------------------------------------------------
\subsection{Prerequisites}
% ------------------------------------------------------------

\begin{statebox}[What You Need Before Starting]
\begin{itemize}
    \item \textbf{KVL and KCL in matrix form} — writing $N$ loop or
          node equations as a single matrix equation.
    \item \textbf{Impedance notation} — $Z_R = R$, $Z_L = pL$,
          $Z_C = \frac{1}{pC}$, where $p = d/dt$ is the Heaviside
          (time-derivative) operator.
    \item \textbf{Matrix multiplication and linear transformations}
          — how multiplying both sides of an equation by a matrix
          changes the basis of the unknowns.
    \item \textbf{Half-bridge converter operation} — the two switch
          states, what the phase voltage $v_A$ looks like, and the
          dc-bus voltage $\Vdc$.
    \item \textbf{What a Fourier series does (conceptually)} — that
          a periodic switched waveform contains harmonics at integer
          multiples of the switching frequency.
\end{itemize}
\end{statebox}

% ------------------------------------------------------------
\subsection{Symbol Table}
% ------------------------------------------------------------

\begin{center}
\renewcommand{\arraystretch}{1.35}
\begin{tabular}{lll}
\toprule
\textbf{Symbol} & \textbf{Meaning} & \textbf{Units} \\
\midrule
$N$                 & Number of lines in the system                  & --- \\
$i_n$               & Current on line $n$                            & A   \\
$v_{nP}$            & Voltage from line $n$ to floating reference $P$ & V   \\
$v_{Pg}$            & Voltage from reference $P$ to ground            & V   \\
$P$                 & Arbitrary floating reference point              & --- \\
$\icm$              & Common-mode current: $\sum_{n=1}^{N} i_n$      & A   \\
$\vcm$              & CM voltage: $\frac{1}{N}\sum_{n=1}^{N} v_{nP}$ & V   \\
$i_{mn}$            & DM current between lines $m$ and $n$           & A   \\
$v_{mn}$            & DM voltage between lines $m$ and $n$           & V   \\
$\Ni$               & DM current normalisation factor (= 2 in paper) & --- \\
$\TiN$, $\TvN$      & Current and voltage transformation matrices     & --- \\
$\iDCM$, $\vDCM$    & Vectors of decomposed DM and CM quantities      & A, V \\
$\mathbf{R}$        & Resistance matrix of $N$-line model             & $\Omega$ \\
$\mathbf{L}$        & Inductance matrix                               & H   \\
$\mathbf{S}$        & Elastance matrix ($\mathbf{S} = \mathbf{C}^{-1}$) & F$^{-1}$ \\
$k$                 & Asymmetry parameter ($< 1$ for real MCPM)      & --- \\
$\Cp$               & Nominal baseplate parasitic capacitance         & F   \\
$\Rh$               & Heat-sink resistance to ground                  & $\Omega$ \\
$k\Cp$              & Source-to-baseplate capacitance (smaller line)  & F   \\
$v_{Q1}$, $v_{Q2}$  & Voltages across upper and lower switches        & V   \\
$\Vdc$              & DC bus voltage                                  & V   \\
CEM                 & Common-mode equivalent model                    & --- \\
MM                  & Mixed-mode (DM and CM entangled)                & --- \\
DCM                 & Decomposed differential-common-mode set         & --- \\
LISN                & Line impedance stabilisation network            & --- \\
MCPM                & Multichip power module                          & --- \\
\bottomrule
\end{tabular}
\end{center}

% ============================================================
% SECTIONS 1–5 — TO BE WRITTEN ONE AT A TIME
% ============================================================

% \section{Defining CM and DM: From Words to Mathematics}
% TODO: §1 — Fig. 1 circuit walk, CM/DM definitions Eqs.(1)-(4),
%            transformation matrices Eqs.(5)-(8), 2-line example

% \section{How Asymmetry Couples the Two Modes}
% TODO: §2 — N-line MM equations (Eq. 9), applying transform (Eq. 10),
%            symmetric decoupled case (Eqs. 14-15), asymmetric S matrix (Eq. 16)

% \section{The Half-Bridge Testbed and the MCPM}
% TODO: §3 — MCPM structure, Fig. 2 walk, LISN explanation,
%            reference A, three-section split (Fig. 3)

% \section{The First CM Equivalent Model}
% TODO: §4 — LISN CM equiv., asymmetric cap CM eq. (Eq. 19),
%            reformulation (Eqs. 20-23), Fig. 4 walk, approx. (Eqs. 24-26),
%            Fig. 5 validation, Tables I-II

% \section{The Alternative CEM: Same Answer, Different Insight}
% TODO: §5 — two-section split (Fig. 6), alternative CEM (Fig. 7),
%            leakage vs. recycling comparison

\end{document}
